\section{Methodology}\label{sec:methods}

\subsection{Case-study dwelling and climate}
The case study is a 95\,m$^2$ top-floor three-bedroom flat in Pune, India
(18.53\,$^\circ$N, 73.85\,$^\circ$E, 559\,m). Three conditioned zones are modelled:
Zone~1, the master bedroom (30\,m$^2$, south and west facades); Zone~2, the living and
dining room (45\,m$^2$, north and west); and Zone~3, the second bedroom (20\,m$^2$,
north and east). The kitchen is unconditioned. The ceiling height is 2.85\,m, so the
conditioned volume is 270.75\,m$^3$.

The envelope consists of:
\begin{itemize}
  \item 200\,mm autoclaved aerated concrete (AAC) walls, $U=0.75$\,W\,m$^{-2}$K$^{-1}$;
  \item a 150\,mm reinforced-concrete roof slab with insulation,
        $U=1.2$\,W\,m$^{-2}$K$^{-1}$ (the Eco-Niwas Samhita limit);
  \item uPVC double glazing, $U=2.8$\,W\,m$^{-2}$K$^{-1}$, SHGC $=0.55$, at a
        window-to-wall ratio of 0.20, with month-dependent overhang (chajja) shading
        of the direct beam.
\end{itemize}
Infiltration is 0.5\,ACH. Table~\ref{tab:params} lists all parameters; every value is
defined in a single configuration file of the released code.

Weather is the ISHRAE typical year for Pune (station 430630). Its dry-bulb temperature
ranges from 4.9 to 41.0\,$^\circ$C, and the hottest hour of the mean day is 13:00--14:00.
The dry pre-monsoon (March--May, mean relative humidity 39--56\,\%) contrasts with the
humid monsoon (June--September, 79--86\,\%). Hourly records are interpolated to the
15-min simulation step: temperatures and humidity ratio linearly, irradiance by holding
each hour's value.

Incident irradiance on each facade is computed from:
\begin{itemize}
  \item the solar position (Spencer declination and equation of time);
  \item EPW direct-normal irradiance on the facade;
  \item isotropic sky-diffuse and ground-reflected components (albedo 0.2).
\end{itemize}
Opaque surfaces see sol-air temperatures. The roof uses absorptance 0.7 and a long-wave
correction of 63\,W\,m$^{-2}$; the walls use absorptance 0.4.

Occupancy follows weekday and weekend schedules for four occupants: bedrooms are
occupied at night, the living room in the morning and evening, and one person is at
home during weekdays. Each occupant adds 75\,W sensible and 55\,W latent heat. Lighting
(5\,W\,m$^{-2}$) and equipment (4.5\,W\,m$^{-2}$) follow fixed profiles. Thirty per
cent of the 2\,kW cooking load enters Zone~2 during the 07:00--09:00 and 19:00--21:00
cooking hours.

\subsection{Thermal model}
Each zone $z\in\{1,2,3\}$ has two nodes. The air node, $T_{a,z}$ (air and furnishings,
capacitance $C_{a,z}$), and the thermal-mass node, $T_{m,z}$ (effective roof and floor
slabs and walls, $C_{m,z}$), are coupled by interior convection $H_{am,z}$:
\begin{align}
C_{a,z}\dot T_{a,z} &= H_{am,z}(T_{m,z}-T_{a,z}) + (UA_{g,z}+H_{inf,z})(T_o-T_{a,z})
  + \sum_{j\neq z} H_{zj}(T_{a,j}-T_{a,z}) \nonumber\\
  &\quad + f_s Q_{sol,z} + f_c Q_{int,z} + Q_{vent,z} - Q_{c,z}, \label{eq:air}\\
C_{m,z}\dot T_{m,z} &= H_{am,z}(T_{a,z}-T_{m,z}) + UA_{w,z}(T_{sa,w,z}-T_{m,z})
  + UA_{r,z}(T_{sa,r}-T_{m,z}) \nonumber\\
  &\quad + (1-f_s)Q_{sol,z} + (1-f_c)Q_{int,z}. \label{eq:mass}
\end{align}
The terms in Eqs.~\eqref{eq:air}--\eqref{eq:mass} are:
\begin{itemize}
  \item $Q_{sol,z}$: transmitted solar gain, of which $f_s=0.3$ goes to the air node;
  \item $Q_{int,z}$: internal gain, convective fraction $f_c=0.6$;
  \item $Q_{vent,z}$: sensible effect of DOAS and IEC supply air;
  \item $Q_{c,z}$: sensible cooling delivered by the indoor unit.
\end{itemize}
In compact form, $\dot{x}=Ax+Bu+Ed$, with states
$x=[T_{a,1..3},T_{m,1..3}]^\top$, inputs $u=Q_{c,1..3}$ and disturbances $d$. The model
is discretised exactly with a zero-order hold at $\Delta t=900$\,s:
\begin{equation}
x_{k+1}=A_d x_k+B_d u_k+E_d d_k,\qquad
\begin{bmatrix}A_d & B_d & E_d\end{bmatrix}
=\begin{bmatrix}I&0&0\end{bmatrix}\exp\!\left(\begin{bmatrix}A&B&E\\0&0&0\end{bmatrix}\Delta t\right).
\end{equation}
Because the discretisation is exact, the update is unconditionally stable and needs no
artificial temperature clamps. Over 24\,h it agrees with fine-step RK4 integration to
$10^{-12}$\,K. The modal time constants are 0.35--0.43\,h for the air nodes and
25--36\,h for the mass nodes (Table~\ref{tab:params}).

\paragraph{Plant--model mismatch.}
The model-based controllers use the nominal parameters, while the simulated "true"
plant deviates from them:
\begin{itemize}
  \item mass capacitance $\times0.85$ and air capacitance $\times1.10$;
  \item envelope conductances $\times1.10$;
  \item solar gains $\times1.10$ and internal gains $\times1.15$.
\end{itemize}
Every controller receives zone air temperatures with Gaussian sensor noise
($\sigma=0.1$\,K).

\subsection{HVAC equipment}
\paragraph{VRF.}
An 11\,kW inverter outdoor unit serves indoor units of 3.5, 5.0 and 2.5\,kW (total
capacity) in Zones~1--3. The model includes:
\begin{itemize}
  \item capacity derating of 0.8\,\%\,K$^{-1}$ above 35\,$^\circ$C;
  \item rated COP 3.8 at 35\,$^\circ$C, with the energy input ratio rising by
        2.5\,\%\,K$^{-1}$;
  \item a part-load curve
        $\mathrm{EIR}(\mathrm{PLR})=0.08+0.62\,\mathrm{PLR}+0.30\,\mathrm{PLR}^2$ down to
        a minimum modulation of 20\,\%, with on--off cycling below it (degradation
        coefficient 0.25);
  \item a sensible heat ratio of 0.75 when the indoor humidity ratio exceeds the coil
        apparatus dew point (0.0087\,kg\,kg$^{-1}$), and 1 otherwise;
  \item 30\,W per running indoor fan;
  \item an operating envelope of 10--50\,$^\circ$C outdoor, outside which the unit
        delivers no cooling and draws no power.
\end{itemize}

\paragraph{Indirect evaporative cooler (IEC).}
The IEC moves 0.35\,m$^3$\,s$^{-1}$ of outdoor air with product-air temperature
$T_{sup}=T_o-0.8(T_o-T_{wb})$ and 180\,W of fan and pump power. One supervisory rule,
identical for all controllers, enables it when four conditions hold: $T_o>28\,^\circ$C,
outdoor RH below 50\,\%, cooling is requested, and
$T_{sup}<\bar T_a-2$\,K. Its sensible effect,
$\dot m c_p (T_{sup}-T_{a,z})$, enters the air nodes. Because the cooler is indirect,
it adds no moisture, but its outdoor air does count in the CO$_2$, PM$_{2.5}$ and
moisture balances.

\paragraph{DOAS.}
Continuous ventilation at the ASHRAE~62.2 whole-dwelling rate (28\,L\,s$^{-1}$) passes
through:
\begin{itemize}
  \item an energy-recovery ventilator (sensible effectiveness 0.70, latent 0.60);
  \item a MERV-13 filter (single-pass PM$_{2.5}$ efficiency 0.6);
  \item a dehumidifying coil (COP 3), which runs only when the indoor humidity ratio
        exceeds 0.0115\,kg\,kg$^{-1}$ and then supplies air at 15\,$^\circ$C and
        $W\le0.0095$.
\end{itemize}
The original study's humidity-triggered DOAS (60\,L\,s$^{-1}$ when outdoor RH $>65$\,\%
and $T_o>22\,^\circ$C) is kept as a scenario.

\subsection{Moisture and indoor air quality}
Indoor humidity ratio $W$, CO$_2$ concentration $C$ and PM$_{2.5}$ concentration $P$ are
well-mixed balances over the conditioned volume $V$:
\begin{align}
\rho V\dot W &= \rho F_{inf}(W_o-W)+\rho F_{d}(W_{d}-W)+\rho F_{iec}(W_o-W)+G_{lat}-\dot m_{coil},\\
V\dot C &= F_{out}(C_o-C)+n_{occ}\,g_{CO_2},\\
V\dot P &= P_o\big(p F_{inf}+(1-\eta_d)F_d+(1-\eta_{iec})F_{iec}\big) + S_{PM}
  -\big(F_{out}+\eta_r F_r + kV\big)P.
\end{align}
Here $F_{out}=F_{inf}+F_d+F_{iec}$ is the exfiltration and exhaust flow balancing the
supply. The remaining parameters are:
\begin{itemize}
  \item $p=0.8$, the infiltration penetration factor;
  \item $k=0.5$\,h$^{-1}$, the deposition rate;
  \item $F_r=0.3$\,m$^3$\,s$^{-1}$ of HEPA H13 recirculation, $\eta_r=0.9995$;
  \item $S_{PM}$, cooking and background emissions;
  \item $g_{CO_2}=5.2$\,mL\,s$^{-1}$ per person.
\end{itemize}
Each balance is linear with coefficients held constant over a step. It is therefore
integrated exactly as $y_{k+1}=y_\infty+(y_k-y_\infty)e^{-\kappa\Delta t}$, which stays
stable for any removal rate $\kappa$. The outdoor PM$_{2.5}$ series is a
Pune-representative monthly--diurnal profile with wet-weather wash-out (Section~\ref{sec:limits}).

\subsection{Control strategies}
All five controllers act on the same plant and the same occupancy-based comfort
schedule. When a zone is occupied, the band is 23--26\,$^\circ$C, the summer range for
air-conditioned buildings in the 2005 National Building Code of India \cite{nbc2005}
(the 2016 edition moved to an adaptive model), and the setpoint is 24.5\,$^\circ$C. When it is unoccupied, the band is 16--29\,$^\circ$C and
the setpoint is 28.5\,$^\circ$C. The original study gave the setpoint setback to only
some controllers; here every controller uses the same schedule.

\paragraph{Thermostat (baseline).}
Each indoor unit runs at full capacity when the measured zone temperature exceeds the
setpoint by 0.5\,K and stops at 0.5\,K below it. This is the behaviour of a typical
residential split air conditioner.

\paragraph{PID.}
Three independent PI loops with a filtered derivative, back-calculation anti-windup and
output limits $[0,\bar u_z]$. Gains follow the SIMC rules
\cite{skogestad2003} for an integrating approximation of the identified zone response,
$K_c=1/\big(k'(\tau_c+\theta)\big)$ and $T_i=\min\big(\tau,4(\tau_c+\theta)\big)$, with
$\tau_c=15$\,min and $\theta=\Delta t/2$.

\paragraph{Economic MPC (E-MPC).}
Every hour the controller solves
\begin{equation}
\min_{U}\;\sum_{h=1}^{N_h}\sum_{z}\frac{\pi_h\,\mathrm{EIR}_T(\hat T_{o,h})}{\mathrm{COP}_r\,\mathrm{SHR}}U_{h,z}
+\rho\sum_{j=1}^{4N_h}\sum_z\Big(\big[\hat T_{a,z,j}-\overline T_{z,j}+\delta\big]_+^2+\big[\underline T_{z,j}+\delta-\hat T_{a,z,j}\big]_+^2\Big)
+\lambda\sum_h\|U_h-U_{h-1}\|^2 \label{eq:empc}
\end{equation}
subject to $0\le U_{h,z}\le\bar u_z$ and $\sum_z U_{h,z}\le \mathrm{SHR}\cdot Q_{cap}(\hat T_{o,h})$.
In Eq.~\eqref{eq:empc}:
\begin{itemize}
  \item the inputs are held for each hour (move blocking) over a horizon of
        $N_h=24$\,h, i.e. 96 model steps;
  \item $\pi_h$ is the ToD tariff;
  \item the comfort constraints are soft, with weight $\rho=50$\,INR\,K$^{-2}$ per step
        and a back-off of $\delta=0.5$\,K;
  \item $\lambda=0.05$\,INR\,kW$^{-2}$ penalises input moves.
\end{itemize}
Predictions $\hat T_a=\hat T_a^{free}+\Gamma U$ start from a steady-state Kalman filter
on the model augmented with three integrating heat-disturbance states, which makes the
predictions offset-free. They use weather forecasts whose dry-bulb error grows as
$0.15\sqrt{h}$\,K (capped at 2\,K) and whose solar error is 10\,\%. The convex problem
is solved with FISTA \cite{beck2009}, using an exact Euclidean projection onto the
box-and-capacity set (120 iterations, warm-started).

\paragraph{Deep Q-network (DQN).}
One parameter-shared double DQN \cite{mnih2015,vanhasselt2016} acts for each zone. Its
15 observable features are:
\begin{itemize}
  \item zone temperature error from the setpoint;
  \item outdoor temperature and the 3-h forecast change;
  \item hour of day and day of year, each as a sine/cosine pair;
  \item tariff, current and next-hour occupancy, and the setpoint;
  \item the previous action;
  \item a one-hot zone index (3 features).
\end{itemize}
It chooses among five capacity levels (0, 25, 50, 75 and 100\,\%) every 15\,min. The
reward is the negative sum of three terms, with the same comfort weight as the E-MPC:
\begin{itemize}
  \item the zone's share of the ToD electricity cost;
  \item the squared comfort-band violation, weighted by $\beta=\rho$;
  \item a switching penalty.
\end{itemize}
The network is a 64--64 ReLU multilayer perceptron trained with Huber loss, Adam
($3\times10^{-4}$), gradient clipping, Polyak target updates and $\epsilon$-greedy
exploration. Training uses 200 fourteen-day episodes. Each episode applies random
perturbations to the typical year: dry-bulb $\pm1.5$\,K, solar gains $\pm15$\,\% and
internal gains $\pm15$\,\%. The checkpoint is selected on greedy validation returns
over five fixed, perturbed fortnights spread across the year, never on exploration
returns. Training and test use the same plant, including the IEC and DOAS.
In both training and deployment the DQN acts through a safety shield, the same kind of
safety layer the Hybrid controller has. The shield:
\begin{itemize}
  \item forbids cooling when a zone is within 0.5\,K of its lower band or more than
        1.5\,K below its setpoint;
  \item enforces at least 50\,\% capacity when an occupied zone is 0.5\,K above its
        upper band.
\end{itemize}
Following shielded reinforcement learning practice \cite{alshiekh2018}, the replay
buffer stores the agent's \emph{proposed} action together with the outcome of the
executed action. A penalty (5\,INR) is added whenever the shield overrides the
proposal, so unsafe proposals learn low values. An unshielded agent is reported as an
ablation.

\paragraph{Hybrid (E-MPC + PI + safety).}
The E-MPC supplies the hourly feed-forward $U_h$ and a 15-min reference trajectory
$\hat T_a$. A PI layer tracks that reference:
\begin{equation}
u_k = \mathrm{sat}\big(U_h + K_h e_k + v_k\big),\qquad e_k = T_{a,k}^{meas}-\hat T_{a,k},
\end{equation}
with half the SIMC gain, leaky back-calculation integration, and activation only while
the MPC is cooling or the zone is within $\delta$ of the upper band. It is intended to reject
intra-hour disturbances that the hourly optimiser does not see: cooking, solar
transients, and IEC/DOAS heat. Its measured contribution is reported in
Section~\ref{sec:sens}. A safety layer enforces $u_z\ge0.8\bar u_z$ when an
occupied zone exceeds the band by 0.5\,K.

\subsection{Experimental design and metrics}
Each controller is simulated over the full typical year at a 15-min step (35{,}040 steps).
The main study is replicated with five seeds. Each seed changes the sensor noise, the
forecast errors and, for the DQN, the training run. Results are reported as the mean
with a 95\,\% $t$-interval. Paired comparisons use daily totals from seed~1, with a
7-day moving-block bootstrap for the annual difference and the Wilcoxon signed-rank
test. Additional single-seed scenarios cover robustness, sensitivity and ablation:
\begin{itemize}
  \item a +3\,K heatwave year;
  \item thermal mass halved and doubled;
  \item a flat tariff;
  \item no plant--model mismatch;
  \item a perfect-information E-MPC, as an upper bound;
  \item MPC horizons of 6 and 12\,h;
  \item DQN action resolutions of 3 and 9 levels;
  \item three ventilation and filtration architectures.
\end{itemize}

The metrics are:
\begin{itemize}
  \item annual electricity by end use;
  \item cost under an hourly three-tier ToD tariff (INR 5.20, 7.15 and 8.95\,kWh$^{-1}$;
        a hypothetical structure kept from the original study, \emph{not} the MSEDCL
        residential tariff). The actual MSEDCL residential structure, a flat energy
        charge with a solar-hours (09:00--17:00) rebate of INR 0.80\,kWh$^{-1}$
        \cite{merc}, is evaluated as a separate scenario;
  \item 15-min peak demand;
  \item the share of VRF energy used in peak-tariff hours;
  \item operational CO$_2$ emissions (0.727\,kg\,kWh$^{-1}$, CEA grid weighted average \cite{cea});
  \item occupied discomfort in K\,h, the sum of overheating above the band and
        \emph{cooling-induced} overcooling below it;
  \item the share of occupied zone-hours with acceptable conditions;
  \item occupied RMSE from the setpoint;
  \item hours with indoor RH above 70\,\%;
  \item CO$_2$ and PM$_{2.5}$ exposure metrics;
  \item compressor starts and controller computation time.
\end{itemize}
Cooling-induced overcooling is measured below $\min(\underline T, T^{free})$, where
$T^{free}$ comes from a free-floating reference run without cooling. Because the
dwelling has no heating, cold hours caused by the weather are not counted against any
controller.

\begin{table}[t]\centering\small
\caption{Main model parameters (full list in \texttt{config/default\_config.m}).}\label{tab:params}
\begin{tabular}{lccc}\toprule
 & Zone 1 & Zone 2 & Zone 3\\\midrule
Floor area (m$^2$) & 30 & 45 & 20\\
Exterior wall length (m) / facades & 11 / S,W & 15 / N,W & 9.2 / N,E\\
Air-node capacitance $C_a$ (MJ\,K$^{-1}$) & 0.55 & 0.83 & 0.37\\
Mass-node capacitance $C_m$ (MJ\,K$^{-1}$) & 10.4 & 15.4 & 7.15\\
Envelope conductance $UA$ (W\,K$^{-1}$) & 86.5 & 124.7 & 63.8\\
Indoor-unit capacity (kW, total) & 3.5 & 5.0 & 2.5\\\midrule
\multicolumn{4}{l}{Modal time constants: 0.35, 0.39, 0.43\,h (air); 24.7, 29.0, 36.0\,h (mass)}\\
\bottomrule\end{tabular}
\end{table}
